% Body of the short summary (summary.tex). Keep consistent with proposal.tex.
\section*{The problem}
The \Ybion{} experiments deliver \SI{369.5}{nm} light through Thorlabs P3-375PM-FC patch cables
(PM-S350-HP: PANDA, pure-silica core, FC/APC). After a fresh polish the coupling efficiency is
excellent. It then falls over days to weeks, faster at higher power, and both the input and the output end
are affected. At \SI{3}{mW} in the \SI{2.4}{\micro m} mode the end faces see
\SI{\approx1.3e5}{W/cm^2}. Because re-polishing restores the throughput, the damage sits at the end faces
(or within micrometres of them), not in the bulk of the fibre. The proposed remedy is to keep the end
faces in an \textbf{oxygen-free atmosphere}, using \textbf{3D-printed cells} with an \textbf{oxygen
absorber} or an inert-gas purge. The MSc project should demonstrate that this halts the degradation, or
find out why it does not.

\section*{What the literature says}
\begin{itemize}
  \item \textbf{Oxygen cannot reach the bulk.} At room temperature \Oxy{} penetrates silica by only nanometres
        per year~\cite{Kajihara2010}, whereas \Htwo{} crosses a fibre in days~\cite{Lemaire1991}. Any oxygen
        effect must act at the end faces, which is where the group's fibres fail.
  \item \textbf{Hydrocarbon contamination: oxygen protects.} Space lasers (Aeolus/ALADIN), UV
        frequency-conversion cavities in ion-trap labs and GaN laser diodes all suffer absorbing carbon
        deposits in nitrogen or vacuum, and are run with oxygen for this
        reason~\cite{Hippler2016,Mondin2017,Hannig2018,Schoedl2005}. Under CW light at an intensity like
        ours, deposits grew fastest at \emph{trace} oxygen levels~\cite{Parthasarathy2024}.
  \item \textbf{Siloxane contamination: oxygen can harm.} Silicones oxidise to non-volatile \ce{SiO_x}.
        At a \SI{410}{nm} fibre end face, \Oxy{} $\geq\SI{100}{ppm}$ \emph{accelerated} such
        deposits~\cite{Fujifilmpatent2006}. This is the strongest support for the group's hypothesis.
  \item \textbf{Oxygen-independent mechanisms.} Particles pulled into the beam, prevented at \SI{313}{nm}
        by a dry \Ntwo{} overpressure~\cite{Marciniak2017}. Laser-induced periodic structures growing out
        of the silica itself, seen at \SI{405}{nm} on both faces of small-core fibres~\cite{Gonschior2013JNCS}.
  \item \textbf{The gap.} Nobody has varied \Oxy{} in a controlled way for near-UV fibre end faces, and there
        are no lifetime data for step-index pure-silica PM fibre at \SI{369.5}{nm}.
\end{itemize}

\section*{Aim and research questions}
Determine whether, and in which atmosphere, the degradation can be stopped, and demonstrate a
low-cost 3D-printed enclosure that does it in routine use. \textbf{RQ1} Where is the loss, and what is it?
\textbf{RQ2} Does removing \Oxy{} change the rate, in which direction, and is there a threshold?
\textbf{RQ3} What are the roles of water, organic and siloxane vapours, and particles?
\textbf{RQ4} Does a sealed absorber cell hold its atmosphere for months and protect a working link?

\section*{Hypotheses and the decisive test}
\begin{center}\footnotesize
\setlength{\tabcolsep}{4pt}
\begin{tabular}{@{}l cccc l@{}}
\toprule
Hypothesis & \arm{AIR} & \arm{CDA} & \arm{N2} & \arm{ABS} & tell-tale sign \\
\midrule
H1 \Oxy-assisted (e.g.\ siloxane $\to$ \ce{SiO_x}) & $\bullet$ & $\bullet$ & $\circ$ & $\circ$ & rate rises with \Oxy; Si/O-rich deposit \\
H2 carbonaceous contamination, \Oxy{} protects & $\bullet$ & $\circ$ & $\bullet/\circ$ & $\bullet/\uparrow$ & C-rich, fluorescent film \\
H3 restructuring of the silica surface & $\bullet$ & $\bullet/\circ$ & $\bullet/\circ$ & $\bullet/\circ$ & ripples along polarisation \\
H4 particles & $\bullet$ & $\circ$ & $\circ$ & $\circ$ & discrete particles \\
H5 colour centres in the glass & $\bullet$ & $\bullet$ & $\bullet$ & $\bullet$ & no surface change \\
H6 mechanical drift & $\bullet$ & $\bullet$ & $\bullet$ & $\bullet$ & recovered by re-peaking \\
\bottomrule
\end{tabular}
\end{center}
{\footnotesize \arm{AIR} lab air; \arm{CDA} clean dry synthetic air; \arm{N2} clean dry nitrogen;
\arm{ABS} sealed printed cell with \Oxy{} absorber. $\bullet$ degrades like air, $\circ$ strongly
reduced, $\uparrow$ faster. \arm{CDA} and \arm{N2} differ \emph{only} in \Oxy, so comparing them
decides the oxygen question.}

\begin{center}
\resizebox{0.66\textwidth}{!}{% Candidate degradation processes at a UV-illuminated fibre end face (input end shown).
\begin{tikzpicture}[font=\sffamily\scriptsize,>=Stealth,x=1cm,y=1cm]
  % --- ferrule, cladding, core
  \fill[gray!22] (0,-1.25) rectangle (6.6,1.25);
  \fill[cyan!7] (0,-0.55) rectangle (6.6,0.55);
  \draw[cyan!45!gray,thin] (0,-0.55) -- (6.6,-0.55) (0,0.55) -- (6.6,0.55);
  \fill[uvviolet!40] (0,-0.11) rectangle (6.6,0.11);
  \draw[gray!65!black,thick] (6.6,-1.25) -- (6.6,1.25);
  \draw[gray!55] (0,-1.25) -- (6.6,-1.25) (0,1.25) -- (6.6,1.25);
  \node[gray!50!black,anchor=west] at (0.1,0.92) {zirconia ferrule};
  \node[gray!50!black,anchor=west] at (0.1,0.36) {silica cladding};
  \node[white,anchor=west] at (0.1,0) {\textbf{core}};
  \draw[decorate,decoration={snake,amplitude=0.6mm,segment length=3mm},gray!60] (-0.02,-1.25) -- (-0.02,1.25);
  % --- converging UV beam
  \fill[uvviolet!22] (6.6,0.10) -- (10.9,1.2) -- (10.9,-1.2) -- (6.6,-0.10) -- cycle;
  \draw[uvviolet!70,thin] (6.6,0.10) -- (10.9,1.2) (6.6,-0.10) -- (10.9,-1.2);
  \draw[->,uvviolet,very thick] (10.75,0) -- (8.9,0);
  \node[uvviolet!80!black,anchor=west] at (9.0,0.24) {\textbf{369.5\,nm}};
  \node[uvviolet!80!black,anchor=west] at (9.0,-0.26) {$\sim\!10^{5}$\,W\,cm$^{-2}$};
  % --- H2: deposit on the core region
  \fill[amberdark!70] (6.6,-0.30) .. controls (6.74,-0.22) and (6.76,0.22) .. (6.6,0.30) -- cycle;
  % --- H5: colour centres in the glass
  \foreach \x in {4.7,5.25,5.7,6.05,6.32}{\node[reddark,inner sep=0] at (\x,0.0) {$\ast$};}
  % --- H1: O2 molecules
  \foreach \x/\y in {7.35/1.55,8.05/1.85,7.0/1.95}{
    \begin{scope}[shift={(\x,\y)}]
      \fill[tealdark] (0,0) circle (0.07); \fill[tealdark] (0.13,0) circle (0.07);
    \end{scope}}
  \draw[->,tealdark,thin] (7.3,1.45) .. controls (6.95,1.1) .. (6.67,0.2);
  % --- H3: adsorbed water
  \foreach \x/\y in {7.25/-1.65,7.9/-1.9}{
    \begin{scope}[shift={(\x,\y)}]
      \fill[cyan!60!blue] (0,0) circle (0.075);
      \fill[gray!60] (-0.09,0.07) circle (0.04); \fill[gray!60] (0.09,0.07) circle (0.04);
    \end{scope}}
  \draw[->,cyan!60!blue,thin] (7.2,-1.55) .. controls (6.9,-1.1) .. (6.66,-0.25);
  % --- H4: particle drawn into the beam
  \fill[black!75] (9.4,-1.55) circle (0.09);
  \draw[->,black!75,thin,densely dashed] (9.3,-1.48) .. controls (8.2,-0.9) .. (6.75,-0.08);
  % --- H6: mechanical drift
  \draw[<->,greydark,thick] (3.2,-1.55) -- (4.4,-1.55);
  % --- annotation boxes
  \tikzset{ann/.style={draw=#1,fill=white,rounded corners=1.5pt,text width=3.55cm,align=left,inner sep=3pt}}
  \node[ann=tealdark,anchor=south west] (h1) at (6.95,2.35)
     {\textbf{\color{tealdark}H1} \Oxy{} reacts with photo-activated surface sites $\to$ absorbing or scattering oxide defects};
  \node[ann=amberdark,anchor=south west] (h2) at (2.95,2.35)
     {\textbf{\color{amberdark}H2} adsorbed organics / siloxanes photo-fixed into a film (oxygen may \emph{protect})};
  \node[ann=reddark,anchor=south west] (h5) at (-1.05,2.35)
     {\textbf{\color{reddark}H5} colour centres in the glass (atmosphere-independent, except \Htwo)};
  \node[ann=cyan!60!blue,anchor=north west] (h3) at (6.95,-2.3)
     {\textbf{\color{cyan!60!blue}H3} adsorbed water assists surface photochemistry or roughening};
  \node[ann=black!75,anchor=north west] (h4) at (2.95,-2.3)
     {\textbf{H4} airborne particles pulled into the high-intensity spot};
  \node[ann=greydark,anchor=north west] (h6) at (-1.05,-2.3)
     {\textbf{\color{greydark}H6} mechanical/thermal drift, epoxy creep (no surface change)};
  \draw[amberdark,thin] (h2.south) -- (6.72,0.27);
  \draw[reddark,thin] (h5.south) -- (5.25,0.09);
  \draw[tealdark,thin] (h1.south) -- (7.55,1.95);
  \draw[cyan!60!blue,thin] (h3.north) -- (7.6,-1.8);
  \draw[black!75,thin] (h4.north) -- (9.35,-1.62);
  \draw[greydark,thin] (h6.north) -- (3.8,-1.62);
\end{tikzpicture}
}
\end{center}

\section*{Approach}
\begin{center}
\resizebox{0.97\textwidth}{!}{% Multi-arm accelerated-ageing testbed: one source, N identical arms, one atmosphere per arm.
\begin{tikzpicture}[font=\sffamily\scriptsize,>=Stealth,x=1cm,y=1cm,
    beam/.style={uvviolet,line width=0.9pt},
    comp/.style={draw=greydark,fill=white,rounded corners=1pt,inner sep=2pt,align=center},
    cell/.style={draw=tealdark,fill=teallight,rounded corners=2pt,thick},
    pd/.style={draw=greydark,fill=gray!35},
    gas/.style={tealdark,line width=0.7pt},
    sig/.style={greydark!80,densely dashed,line width=0.4pt}]
  % ------------------------------------------------------------ source and splitter column
  \node[comp,fill=uvlight,minimum width=2.1cm,minimum height=0.75cm] (laser) at (1.2,8.25)
       {\textbf{laser}\\375\,nm diode or 369.5\,nm};
  \node[comp,minimum width=2.1cm] (shape) at (1.2,7.15) {beam shaping (+ isolator)};
  \draw[beam,->] (laser.south) -- (shape.north);
  \draw[beam] (shape.south) -- (1.2,1.5);
  \foreach \y/\lab/\desc/\celltype in {6.0/1/{ambient lab air (open)}/0,
                                       4.5/2/{clean dry synthetic air, purged}/1,
                                       3.0/3/{clean dry \Ntwo{}, purged}/1,
                                       1.5/4/{sealed cell + \Oxy{} absorber}/2}{
    % splitter: lambda/2 plate + PBS (last arm: mirror)
    \ifnum\lab<4
      \draw[greydark,fill=uvlight] (1.2-0.22,\y-0.22) rectangle (1.2+0.22,\y+0.22);
      \draw[greydark] (1.2-0.22,\y-0.22) -- (1.2+0.22,\y+0.22);
      \draw[greydark,fill=white] (0.95,\y+0.42) rectangle (1.45,\y+0.48);
    \else
      \draw[greydark,line width=1.6pt] (1.2-0.2,\y-0.2) -- (1.2+0.2,\y+0.2);
    \fi
    \draw[beam,->] (1.42,\y) -- (2.75,\y);
    % pick-off wedge with reference / back-reflection photodiodes
    \draw[greydark,fill=cyan!10] (2.85,\y-0.32) -- (3.0,\y-0.32) -- (3.1,\y+0.32) -- (2.95,\y+0.32) -- cycle;
    \draw[beam,densely dotted] (2.98,\y) -- (2.98,\y+0.55);
    \fill[pd] (2.80,\y+0.55) rectangle (3.16,\y+0.66);
    \draw[beam,densely dotted] (3.0,\y) -- (3.0,\y-0.55);
    \fill[pd] (2.82,\y-0.66) rectangle (3.18,\y-0.55);
    \draw[beam,->] (3.1,\y) -- (3.95,\y);
    % focusing asphere
    \draw[greydark,fill=cyan!10] (4.15,\y) ellipse (0.08 and 0.33);
    % input-end atmosphere
    \ifnum\celltype=0
      \draw[greydark!60,densely dashed,rounded corners=2pt] (4.55,\y-0.42) rectangle (5.95,\y+0.42);
    \else
      \draw[cell] (4.55,\y-0.42) rectangle (5.95,\y+0.42);
      \fill[cyan!25] (4.62,\y-0.3) rectangle (4.72,\y+0.3);
    \fi
    \fill[uvviolet!25] (4.23,\y+0.22) -- (5.5,\y+0.03) -- (5.5,\y-0.03) -- (4.23,\y-0.22) -- cycle;
    \fill[gray!55] (5.5,\y-0.09) rectangle (5.95,\y+0.09);
    % fibre (coil) between the two cells
    \draw[uvviolet!70!black,line width=0.8pt] (5.95,\y) -- (6.45,\y)
      .. controls (6.75,\y+0.38) and (7.15,\y+0.38) .. (7.3,\y)
      .. controls (7.45,\y-0.38) and (7.85,\y-0.38) .. (8.0,\y)
      .. controls (8.15,\y+0.38) and (8.55,\y+0.38) .. (8.7,\y) -- (9.2,\y);
    % output-end atmosphere
    \ifnum\celltype=0
      \draw[greydark!60,densely dashed,rounded corners=2pt] (9.2,\y-0.42) rectangle (10.6,\y+0.42);
    \else
      \draw[cell] (9.2,\y-0.42) rectangle (10.6,\y+0.42);
      \fill[cyan!25] (10.43,\y-0.3) rectangle (10.53,\y+0.3);
    \fi
    \fill[gray!55] (9.2,\y-0.09) rectangle (9.65,\y+0.09);
    \fill[uvviolet!25] (9.65,\y+0.03) -- (11.25,\y+0.3) -- (11.25,\y-0.3) -- (9.65,\y-0.03) -- cycle;
    % output detector (integrating sphere)
    \draw[pd] (11.6,\y) circle (0.3);
    \fill[white] (11.3,\y-0.07) rectangle (11.36,\y+0.07);
    % label
    \node[anchor=west,align=left] at (12.05,\y) {\textbf{arm \lab}\\\desc};
    \ifnum\celltype=2
      \fill[amberdark!75] (5.15,\y-0.40) rectangle (5.45,\y-0.28);
      \fill[orange] (5.25,\y+0.33) circle (0.05);
      \fill[amberdark!75] (9.85,\y-0.40) rectangle (10.15,\y-0.28);
      \fill[orange] (9.95,\y+0.33) circle (0.05);
    \fi
  }
  % ------------------------------------------------------------ annotations of the first arm
  \node[anchor=south,align=center] at (4.15,6.45) {asphere\\$f\!\approx\!5$\,mm};
  \node[anchor=south,align=center] at (5.25,6.45) {input\\end};
  \node[anchor=south] at (7.6,6.45) {1--2\,m PM fibre};
  \node[anchor=south,align=center] at (9.9,6.45) {output\\end};
  \node[anchor=south,align=center] at (11.6,6.45) {$P_\text{out}$};
  \node[anchor=north,align=center] at (3.0,0.78) {pick-off wedge:\\$P_\text{ref}$ (up), $P_\text{back}$ (down)};
  % purged cells: gas in / out stubs
  \foreach \y in {4.5,3.0}{
    \draw[gas,->] (5.05,\y+0.68) -- (5.05,\y+0.42);
    \draw[gas,->] (5.45,\y+0.42) -- (5.45,\y+0.68);
    \draw[gas,->] (9.7,\y+0.68) -- (9.7,\y+0.42);
    \draw[gas,->] (10.1,\y+0.42) -- (10.1,\y+0.68);
  }
  % ------------------------------------------------------------ gas panel
  \node[comp,fill=teallight,draw=tealdark,minimum width=2.4cm,align=left] (cyl) at (1.75,-0.95)
       {\Ntwo{} (5.0 or LN$_2$ boil-off)\\synthetic air (zero grade)\\opt.: forming gas, \Oxy/\Ntwo{} mix};
  \node[comp,draw=tealdark,align=center] (pur) at (5.0,-0.95) {regulators, purifier,\\3\,nm particle filter};
  \node[comp,draw=tealdark,align=center] (flow) at (7.85,-0.95) {flow control\\5--20\,sccm per cell};
  \draw[gas,->] (cyl.east) -- (pur.west);
  \draw[gas,->] (pur.east) -- (flow.west);
  \draw[gas,->] (flow.east) -- (9.35,-0.95) node[anchor=west,align=left,tealdark] {to the cells\\of arms 2 and 3};
  \node[tealdark,anchor=north west,align=left] at (0.55,-1.65)
       {cell outlets $\to$ RH sensor, trace \Oxy{} analyser, VOC sampling $\to$ long exhaust line or check valve (no oil bubbler)};
  % ------------------------------------------------------------ data acquisition
  \node[comp,minimum width=2.4cm,align=left,anchor=west] (daq) at (12.0,-0.95)
       {\textbf{DAQ + PC}\\photodiodes, \Oxy, RH, $T$, $p$,\\flows, laser current};
  % ------------------------------------------------------------ legend
  \begin{scope}[shift={(0.55,-2.45)}]
    \draw[cell] (0,0) rectangle (0.35,0.22); \node[anchor=west] at (0.4,0.11) {end-face cell (3D-printed, UV window)};
    \fill[amberdark!75] (5.1,0.04) rectangle (5.4,0.18); \node[anchor=west] at (5.45,0.11) {\Oxy{} absorber};
    \fill[orange] (7.5,0.11) circle (0.05); \node[anchor=west] at (7.6,0.11) {\Oxy{} sensor spot};
    \fill[pd] (9.75,0.04) rectangle (10.05,0.18); \node[anchor=west] at (10.1,0.11) {photodiode};
    \draw[gas] (11.7,0.11) -- (12.1,0.11); \node[anchor=west] at (12.15,0.11) {clean gas};
  \end{scope}
\end{tikzpicture}
}
\end{center}
\begin{itemize}
  \item \textbf{Phase 0, forensics (months 1--2).} Image, test and analyse fibres already degraded in the
        lab (microscopy, SEM/EDX, AFM, Raman), and survey contaminant sources such as silicones, printed
        parts and cable jackets.
  \item \textbf{Testbed.} A dedicated \SI{375}{nm} single-mode diode (\SI{70}{mW}) feeds 4--6
        identical arms in \emph{parallel}. For each arm the testbed logs throughput, end-face reflection,
        \Oxy, RH, VOC index, temperature and flow continuously. A sentinel arm tracks system drift, and a
        daily re-peak separates alignment drift from end-face change.
  \item \textbf{End-face cells.} 3D-printed (sealed and baked) with a UV window, the connector fully
        enclosed, a potted cable gland, gas swept across the face, and an absorber behind a PTFE membrane.
        An optical \Oxy{} sensor spot reads the cell from outside. Target: $<\SI{100}{ppm}$ \Oxy{} for months,
        leak rate $<\SI{1e-5}{mbar.L/s}$.
  \item \textbf{Campaigns.} Baseline in air at 1, 3 and \SI{10}{mW} (power law). Campaign I: \arm{AIR},
        \arm{CDA}, \arm{N2}, \arm{ABS}. Campaign II: replicates, \Oxy{} dose--response (10~ppm to
        \SI{1}{\percent}), humid gas, polymer and silicone challenges, crossed input/output. Finally,
        demonstration on a working \SI{369.5}{nm} link for $\geq8$ weeks.
  \item \textbf{Locating the loss.} At the end of each run, a fixed cleaning sequence (dry, solvent,
        \Oxy{} plasma or UV/ozone, re-polish) is applied to each face in turn, with a measurement after
        every step. Uncleaned twins go to surface analysis.
\end{itemize}

\section*{Work plan (12 months; a 6-month variant exists)}
\begin{center}
\resizebox{0.9\textwidth}{!}{% 12-month work plan in half-month units (1 unit = ~2 weeks).
\begin{ganttchart}[
    x unit=0.43cm, y unit chart=0.5cm, y unit title=0.5cm,
    hgrid={*1{draw=gray!20}}, vgrid={*1{draw=gray!12},*1{draw=gray!35}},
    time slot format=simple, canvas/.append style={draw=none},
    title/.append style={fill=uvviolet!12,draw=gray!40}, title label font=\sffamily\scriptsize\bfseries,
    bar/.append style={fill=uvviolet!55,draw=uvviolet!80!black}, bar height=0.55,
    bar label font=\sffamily\scriptsize, bar label node/.append style={align=right},
    group/.append style={fill=tealdark,draw=tealdark}, group label font=\sffamily\scriptsize\bfseries,
    group right shift=0, group top shift=0.35, group height=0.3, group peaks height=0.15,
    milestone/.append style={fill=amberdark,draw=amberdark!70!black}, milestone height=0.55,
    milestone label font=\sffamily\scriptsize\itshape\color{amberdark!80!black},
    milestone label node/.append style={align=right},
    link/.style={-Stealth,draw=gray!60}
  ]{1}{24}
  \gantttitle{Month}{24} \\
  \gantttitlelist{1,...,12}{2} \\
  \ganttgroup{Set-up}{1}{10} \\
  \ganttbar{WP0 Induction, safety, reading}{1}{3} \\
  \ganttbar{WP1 Forensics of used fibres}{2}{4} \\
  \ganttbar{WP2 Testbed + data logging}{3}{8} \\
  \ganttbar{WP3 End-face cells: design, print, qualify}{4}{10} \\
  \ganttmilestone{MS1 Plan + risk assessments signed off}{2} \\
  \ganttmilestone{MS2 Forensic report; MS3 testbed qualified}{8} \\
  \ganttgroup{Experiments}{7}{22} \\
  \ganttbar{WP4 Baseline in air: power + dose law}{7}{10} \\
  \ganttbar{WP5 Campaign I: atmosphere screening}{11}{14} \\
  \ganttbar{WP6 Campaign II: confirmation + dose-response}{15}{18} \\
  \ganttbar{WP7 Ex-situ surface analysis}{12}{20} \\
  \ganttbar{WP8 Demonstration on the 369.5\,nm system}{17}{22} \\
  \ganttmilestone{MS4 Cells qualified ($<\!100$\,ppm \Oxy{} for 2 weeks)}{10} \\
  \ganttmilestone{MS5 Campaign I decision}{14} \\
  \ganttmilestone{MS6 Campaign II result}{18} \\
  \ganttgroup{Write-up}{8}{24} \\
  \ganttbar[bar/.append style={fill=uvviolet!25}]{WP9 Methods chapter (while building)}{8}{12} \\
  \ganttbar{WP9 Analysis, thesis, paper draft}{18}{24} \\
  \ganttmilestone{MS7 Thesis submitted}{24}
\end{ganttchart}
}
\end{center}

\section*{Success criteria and deliverables}
\begin{itemize}
  \item The effect of \Oxy{} is measured with a confidence interval: a ratio of rates between \arm{N2} and
        \arm{CDA} different from 1 by a factor of at least 2 in two runs, or excluded.
  \item ``Halted'' means a 95\,\% upper limit below \SI{1}{\percent} per \SI{1000}{h} at \SI{3}{mW} (less than
        \SI{10}{\percent} per year of continuous use), at least 20 times below the air arm, with unchanged end
        faces. This is reachable in 3--6 weeks at \SI{0.3}{\percent} daily scatter.
  \item A qualified cell design (CAD, bill of materials, SOPs) tested on a working link; data, code,
        thesis, and a short paper if the results allow.
\end{itemize}

\section*{Resources, risks and first steps}
\textbf{Resources.} The main new items are the \SI{375}{nm} diode with mount and driver, 8--12
P3-375PM-FC cables, windows, UV-resistant photodiodes and DAQ, gas purification and flow control, and
optical \Oxy{} sensing. Much of the optics, the power meters, the \Ntwo{} supply, the 3D printers and the
polishing kit are likely to be in the group already. Indicative costs are in the full proposal.
\textbf{Main risks.} The testbed may not reproduce the lab degradation (mitigated by the forensics first
and by tuning the power). Alignment drift may mask the effect (re-peak protocol, reflection monitor). The
cells or the printed material itself may contaminate the faces (leak tests, material screening, the
\arm{N2+P} arm). Polishing scatter may hide small effects (one procedure, replicates).
\textbf{First month.} Inductions and risk assessments; the five key papers; hands-on cleaning,
inspection and coupling; collecting degraded fibres; a project plan and order list presented to the
group.
